\documentclass[11pt]{article}

\usepackage[margin=1in]{geometry}
\usepackage{amsmath,amssymb,amsthm,mathtools}
\usepackage{graphicx}
\usepackage{booktabs}
\usepackage{enumitem}
\usepackage{xcolor}
\usepackage{hyperref}
\usepackage{microtype}
\usepackage{caption}
\usepackage{subcaption}

\hypersetup{
  colorlinks=true,
  linkcolor=blue!60!black,
  citecolor=blue!60!black,
  urlcolor=blue!60!black
}

\newtheorem{theorem}{Theorem}
\newtheorem{proposition}{Proposition}
\newtheorem{lemma}{Lemma}
\newtheorem{definition}{Definition}
\newtheorem{remark}{Remark}

\newcommand{\R}{\mathbb{R}}
\newcommand{\Tr}{\operatorname{Tr}}
\newcommand{\diag}{\operatorname{diag}}
\newcommand{\eig}{\operatorname{eig}}
\newcommand{\vecop}{\operatorname{vec}}
\newcommand{\relu}{\operatorname{ReLU}}
\newcommand{\Sym}{\mathfrak{S}}
\newcommand{\MLP}{\mathrm{MLP}}
\newcommand{\Spec}{\mathcal{S}_{\mathrm{spec}}}
\newcommand{\TraceShadow}{\mathcal{T}}
\newcommand{\Canon}{\mathcal{C}}
\newcommand{\Had}{\odot}

\title{\vspace{-1.5em}\textbf{Weight Shadows: Quotienting Transformer Parameter Space by Native Symmetries}}
\author{Vinicius Santos}
\date{Exploratory draft, \today}

\begin{document}
\maketitle

\begin{abstract}
Neural networks are usually compared as points in weight space. But weight space contains exact algebraic redundancies: permuting hidden units, positively rescaling ReLU channels, or reordering attention heads can move a model far in Euclidean parameter distance while leaving the computed function unchanged. This note introduces \emph{weight shadows}: algebraic invariants of transformer blocks that descend from parameter representatives to symmetry orbits. For an MLP sublayer $x\mapsto F_2\relu(F_1x)$, we attach the coupling matrix
\[
G=(F_1F_1^T)\Had(F_2^TF_2),
\]
whose entries pair incoming and outgoing hidden-unit interactions. The native invariant is the trace-power shadow $(\Tr G,\Tr G^2,\ldots)$; the sorted eigenvalue multiset is its stable spectral readout. This trace/spectral shadow is sound under hidden-unit permutation and positive ReLU scaling, but incomplete because it is blind to the difference between permutation conjugacy and general orthogonal similarity. We then introduce a stronger Level-4 canonical coupling shadow that retains node-localized cross-terms and recovers a generic canonical form for the coupling graph. On trained microGPT blocks, raw weight distance ranks function-identical symmetry copies as farther from the original than real continued-training fine-tunes, while the shadow ranks them correctly. A model-merging test further marks the scope boundary: canonical shadows identify orbits, but do not by themselves transport between distinct orbits. The result is a concrete transformer-domain instance of the broader principle: do not compare representatives when the native object is an orbit; compare the shadow that descends to the quotient.
\end{abstract}

\paragraph{Keywords.} transformer blocks, neural network symmetries, weight space, quotient invariants, graph canonization, checkpoint deduplication, shadow algebra.

\section{The paradox: same function, farther weights}

Suppose two transformer blocks compute the same function. Not approximately the same. Exactly the same, up to floating-point roundoff. One block is obtained from the other by merely reordering hidden units and compensating the corresponding output columns. Another block is obtained by continuing training for a few optimization steps, so its function genuinely changes. Which block should be closer to the original?

Euclidean weight distance gives the wrong answer. In our microGPT experiment, a symmetry copy of a trained block sits at weight distance about $7.1$ from the original while its block function is unchanged. Real fine-tunes sit at weight distance about $1.1$--$1.2$ while the function changes substantially. Raw parameter distance therefore ranks an identical function as farther away than a changed function.

This is not a numerical oddity. It is a category error. Weight distance compares \emph{representatives}. The block function is constant on a \emph{symmetry orbit}. When the meaningful object lives on a quotient, the first task is not to find a better norm on representatives. It is to construct an invariant that descends to the quotient.

\begin{figure}[t]
  \centering
  \includegraphics[width=0.98\linewidth]{figures/sweep_result.png}
  \caption{The basic pathology and the shadow correction. Left: a function-identical symmetry copy can be far in raw weight distance, while perturbations and real fine-tunes can be much nearer in weight space despite changing the function. Right: in the weight-shadow plane, the symmetry copy collapses to the same-function band while real fine-tunes separate.}
  \label{fig:paradox}
\end{figure}

The figure is the whole motivation. The star is a symmetry copy: far in weights, zero in function distance. The orange squares are real continued-training fine-tunes: near in weights, changed in function and shadow. The plot says what the algebra will formalize:
\[
\text{weights are representatives;}\qquad \text{functions live on quotients;}\qquad \text{shadows measure the quotient.}
\]

\section{From phase shadows to weight shadows}

The construction follows a pattern from the phase-shadow program. In the Fibonacci anyon setting, braid matrices carry a nuisance symmetry: global phase. The meaningful matrix is not a representative $M$ but its projective class $\lambda M$. A phase-correlation shadow removes that nuisance by using products of the form $M_a\overline{M_b}$ rather than comparing entries directly.

Transformer blocks have a different nuisance group. For an MLP sublayer, hidden units can be permuted. For ReLU networks, hidden units can also be positively rescaled: multiply an incoming row by $a>0$ and divide the corresponding outgoing column by $a$. For multi-head attention, heads can be reordered when the output projection is reordered accordingly. These transformations move the parameter tensor, but they do not change the represented function.

Thus the transfer is direct:
\[
\begin{array}{ccc}
\text{anyon matrices} & \longrightarrow & \text{transformer weights}\\[2pt]
\text{global phase} & \longrightarrow & \text{permutation and scaling gauges}\\[2pt]
\text{phase-correlation shadow} & \longrightarrow & \text{weight-coupling shadow}.
\end{array}
\]

The goal is not to make neural networks ``golden'' by decoration. The goal is to use the same native algebraic discipline: identify the primitive nuisance symmetry, build the coupling that cancels it, and only then measure distance.

\section{The native symmetry group of a ReLU MLP block}

Consider a one-hidden-layer ReLU sublayer
\begin{equation}
  f_{F_1,F_2}(x)=F_2\relu(F_1x),
  \label{eq:mlp}
\end{equation}
where $F_1\in\R^{H\times d}$ and $F_2\in\R^{d'\times H}$. Write
\[
 r_i=F_1[i,:]\in\R^d,\qquad c_i=F_2[:,i]\in\R^{d'}.
\]
The contribution of hidden unit $i$ is $c_i\relu(r_ix)$.

There are two exact symmetries.

\paragraph{Permutation.} For any permutation matrix $P\in\R^{H\times H}$,
\[
F_1' = PF_1,\qquad F_2'=F_2P^T
\]
gives
\[
F_2'\relu(F_1'x)=F_2P^T\relu(PF_1x)=F_2P^TP\relu(F_1x)=F_2\relu(F_1x).
\]

\paragraph{Positive ReLU scaling.} For any positive diagonal matrix $D=\diag(a_1,\ldots,a_H)$ with $a_i>0$,
\[
F_1'=DF_1,\qquad F_2'=F_2D^{-1}
\]
gives
\[
F_2'\relu(F_1'x)=F_2D^{-1}\relu(DF_1x)=F_2D^{-1}D\relu(F_1x)=F_2\relu(F_1x).
\]

Together these generate the MLP gauge group
\begin{equation}
  \Gamma_{\MLP}=(\R_{>0})^H\rtimes \Sym_H.
  \label{eq:gauge-group}
\end{equation}
The weights $(F_1,F_2)$ are representatives. The function $f_{F_1,F_2}$ is constant on $\Gamma_{\MLP}$-orbits.

\section{The coupling matrix}

A single hidden-unit descriptor $(r_i,c_i)$ is not invariant under positive scaling. But the pairwise incoming-outgoing coupling
\begin{equation}
  G_{ij}=\langle r_i,r_j\rangle\,\langle c_i,c_j\rangle
  \label{eq:coupling-entry}
\end{equation}
is invariant under that scaling. In matrix form,
\begin{equation}
  G=(F_1F_1^T)\Had(F_2^TF_2),
  \label{eq:coupling}
\end{equation}
where $\Had$ denotes the Hadamard product.

\begin{proposition}[Scaling cancellation]
The coupling matrix $G$ is invariant under positive ReLU rescaling of hidden units.
\end{proposition}

\begin{proof}
Under $r_i\mapsto a_ir_i$ and $c_i\mapsto a_i^{-1}c_i$,
\[
\langle r_i,r_j\rangle\mapsto a_ia_j\langle r_i,r_j\rangle,
\]
while
\[
\langle c_i,c_j\rangle\mapsto a_i^{-1}a_j^{-1}\langle c_i,c_j\rangle.
\]
The factors cancel in the product \eqref{eq:coupling-entry}, so every entry $G_{ij}$ is unchanged.
\end{proof}

\begin{proposition}[Permutation equivariance]
Under hidden-unit permutation $F_1\mapsto PF_1$, $F_2\mapsto F_2P^T$, the coupling matrix transforms as
\[
G\mapsto PGP^T.
\]
\end{proposition}

\begin{proof}
The incoming Gram matrix transforms as $F_1F_1^T\mapsto P(F_1F_1^T)P^T$, and the outgoing Gram matrix transforms as $F_2^TF_2\mapsto P(F_2^TF_2)P^T$. Their Hadamard product is therefore $P[(F_1F_1^T)\Had(F_2^TF_2)]P^T$.
\end{proof}

Thus $G$ has already done the first quotient step. It cancels the continuous scaling gauge exactly and leaves only a finite relabeling ambiguity.

\section{The trace shadow and its spectral readout}

The coupling matrix $G$ is scaling-invariant and permutation-equivariant. To remove the remaining hidden-unit ordering, we take conjugation-invariant signatures.

The framework-native choice is the sequence of trace powers:
\begin{equation}
  \TraceShadow(G)=\bigl(\Tr(G),\Tr(G^2),\ldots,\Tr(G^H)\bigr).
  \label{eq:trace-shadow}
\end{equation}
We call this the \emph{trace shadow}.

This definition uses multiplication, addition, and trace. It does not require diagonalization, eigenvectors, roots, or a spectral completion. Expanding the $k$th component gives
\begin{equation}
\Tr(G^k)=\sum_{i_1,\ldots,i_k}G_{i_1i_2}G_{i_2i_3}\cdots G_{i_ki_1}.
\label{eq:closed-walks}
\end{equation}
So the trace shadow measures closed cycles in the hidden-unit coupling graph. It does not name the hidden units. It counts the interaction cycles that survive the quotient by permutation.

\begin{theorem}[Soundness of the trace shadow]
If two MLP blocks are related by hidden-unit permutation and positive ReLU scaling, then they have the same trace shadow.
\end{theorem}

\begin{proof}
Positive scaling leaves $G$ unchanged. A hidden-unit permutation sends $G$ to $PGP^T$. For every $k\ge 1$,
\[
\Tr\bigl((PGP^T)^k\bigr)=\Tr(PG^kP^T)=\Tr(G^k).
\]
Therefore all trace powers agree.
\end{proof}

For numerical comparison, however, raw trace powers are poorly conditioned. The higher powers $\Tr(G^k)$ can span many orders of magnitude, and direct Euclidean comparison of trace vectors becomes dominated by the largest powers. The trace shadow is the conceptually native invariant; it is not the best coordinate system for measuring distances.

The computational readout used in the experiments is therefore the sorted spectrum
\begin{equation}
  \Spec(G)=\operatorname{sort}(\eig(G)).
  \label{eq:spectral-shadow}
\end{equation}
This is not a different invariant in finite dimension. If $\lambda_1,\ldots,\lambda_H$ are the eigenvalues of $G$, then
\[
\Tr(G^k)=\sum_{i=1}^H\lambda_i^k.
\]
Conversely, the first $H$ power sums determine the characteristic polynomial by Newton identities, and hence determine the eigenvalue multiset over the spectral completion.

The hierarchy is:
\[
G\quad\longrightarrow\quad \bigl(\Tr G,\Tr G^2,\ldots,\Tr G^H\bigr)\quad\longleftrightarrow\quad \operatorname{sort}(\eig(G)).
\]
The middle object is the native shadow. The final object is the stable numerical readout.

\section{Why trace/spectral shadows are not complete}

Soundness is not completeness. The trace/spectral shadow descends correctly to the permutation quotient, but it forgets too much: it is invariant under arbitrary similarity, in particular under orthogonal conjugation
\[
G\mapsto QGQ^T,
\]
while the true hidden-unit symmetry only allows permutation conjugation
\[
G\mapsto PGP^T.
\]
Since permutation matrices form a strict subset of the orthogonal group, the trace/spectral shadow is blind to distinctions it should keep if the goal is a canonical form for the permutation quotient.

Equivalently, $G$ is a weighted graph on hidden units. The spectrum of a graph is a classic incomplete invariant: cospectral non-isomorphic graphs exist. Our setting is the weighted version of the same phenomenon. The spectrum gives a fast and useful rejection oracle:
\[
\Spec(G)\neq\Spec(G')\quad\Longrightarrow\quad G'\neq PGP^T\text{ for all permutations }P.
\]
But equal spectrum does not imply permutation equivalence.

This limitation is not a failure of the shadow program. It identifies the next native layer. The diagonal/spectral terms are not enough; the shadow must retain cross-terms tied to the standard hidden-unit basis.

\section{Level 4: the canonical coupling shadow}

The spectral shadow fails because it throws away which neuron is which. A permutation merely relabels the standard basis. A general rotation destroys it. The fix is to retain node-localized information before quotienting by order.

For each hidden unit $i$, define a node signature
\begin{equation}
  \nu_i(G)=\left(G_{ii},\;\sum_jG_{ij},\;\sum_jG_{ij}^2\right).
  \label{eq:node-signature}
\end{equation}
These are self-weight, weighted degree, and incident energy. A permutation reorders the signatures. A generic orthogonal rotation changes them.

On the open stratum where the node signatures are pairwise distinct, they define a unique canonical order. Let $\pi_G$ be the order obtained by sorting the signatures lexicographically. The Level-4 canonical coupling shadow is
\begin{equation}
  \Canon(G)=\vecop\bigl(G[\pi_G,\pi_G]\bigr).
  \label{eq:canon-shadow}
\end{equation}

\begin{theorem}[Generic completeness for coupling-graph permutation]
Assume $G$ and $G'$ have pairwise distinct node signatures. Then
\[
\Canon(G)=\Canon(G')
\]
if and only if $G'$ is permutation-conjugate to $G$.
\end{theorem}

\begin{proof}
If $G'=PGP^T$, then the node signatures of $G'$ are the permuted node signatures of $G$. Since the signatures are distinct, sorting them recovers the same canonical node order, so the reordered matrices agree.

Conversely, if the canonical reordered matrices agree, then undoing the two sorting permutations gives $G'=PGP^T$ for the corresponding permutation $P$.
\end{proof}

This is the exact analogue of moving from diagonal data to cross-correlation data in the phase-shadow construction. The spectral shadow keeps only global class functions of $G$. The canonical coupling shadow keeps enough node-localized structure to distinguish a permutation from a rotation on the generic stratum. This is completeness for the coupling graph $G$ modulo permutation, not for the full functional equivalence class of the MLP.

The hard residual is also clear. If signatures tie, then naive sorting is not a canonical form. The general problem becomes weighted graph canonization: choose a canonical representative of $G$ modulo $G\sim PGP^T$. Weisfeiler--Lehman refinement and exact graph-canonization algorithms extend the Level-4 construction to larger strata, but highly symmetric automorphic cases remain the natural hard set. Trained dense weights generically avoid this set; symmetric toy constructions need explicit tie handling.

\section{Attention-head shadows}

The same idea applies to standard concatenated multi-head attention. In the convention used here, each head has its own rows in $W_Q,W_K,W_V$, the head outputs are concatenated, and the corresponding block of columns in $W_O$ maps the concatenation back to the residual stream. Reordering heads together with the matching $W_O$ column blocks preserves the function.

For head $h$, collect the descriptor
\[
u_h=\vecop(W_Q^{(h)},W_K^{(h)},W_V^{(h)},W_O^{(h)}).
\]
A simple head coupling matrix is
\[
H_{ab}=\langle u_a,u_b\rangle.
\]
The head trace/spectral shadow is then defined exactly as above. This gives a block shadow by concatenating MLP and head shadows:
\[
\mathcal S_{\mathrm{block}}=\mathcal S_{\MLP}\oplus\mathcal S_{\mathrm{head}}.
\]
For the experiments below, this simple head shadow is sufficient to test invariance under head permutation. A richer Level-4 head shadow would canonicalize the head coupling matrix in the same way as the MLP coupling matrix.

\section{Experiments}

The experiments use a minimal microGPT-style character model: embedding dimension $16$, block size $16$, four attention heads, an MLP width of $4d$, and trained weights from multiple random seeds. The point is not scale. The point is that every operation is visible: MLP rows and columns, attention heads, residuals, normalizations, and output projections are explicit.

We compare three distances:
\begin{align*}
 d_W(\theta,\theta')&=\|\theta-\theta'\|_2,\\
 d_f(\theta,\theta')&=\max_X\|f_\theta(X)-f_{\theta'}(X)\|_\infty,\\
 d_{\mathcal S}(\theta,\theta')&=\|\mathcal S(\theta)-\mathcal S(\theta')\|_2.
\end{align*}
Here $d_W$ is representative-level weight distance, $d_f$ is block-function distance on held-out random inputs, and $d_{\mathcal S}$ is shadow distance.

\subsection{Soundness under exact symmetries}

Starting from a trained model, we apply MLP permutations, attention-head permutations, and MLP permutation plus positive scaling. These transformations preserve the block function. They also move the raw weights substantially. The shadow collapses them back to the same point.

\begin{table}[h]
\centering
\caption{Symmetry transformations preserve function and shadow while moving raw weights. Values are representative outputs from the microGPT block experiment.}
\label{tab:soundness}
\begin{tabular}{lccc}
\toprule
Transformation & function distance & weight distance & shadow distance\\
\midrule
MLP permutation & $\approx 10^{-15}$ & $6.41$ & $8\times 10^{-16}$\\
Head permutation & $\approx 10^{-15}$ & $5.39$ & $2\times 10^{-15}$\\
Permutation + scaling & $\approx 10^{-15}$ & $6.92$ & $1\times 10^{-15}$\\
\bottomrule
\end{tabular}
\end{table}

The result is exactly the soundness theorem in numerical form: the representative moves, the orbit does not.

\subsection{Weight distance can rank functional identity backwards}

The anti-correlation concern is not an artifact of one hand-picked noise scale. We swept perturbation magnitude $\sigma$ and compared each perturbation to a same-function symmetry copy at weight distance about $7.1$.

\begin{table}[h]
\centering
\caption{Perturbation sweep. For $\sigma\leq 0.10$, changed models are closer in raw weight distance than a function-identical symmetry copy.}
\label{tab:sweep}
\begin{tabular}{cccc}
\toprule
$\sigma$ & function distance & weight distance & closer than symmetry copy?\\
\midrule
$0.005$ & $0.08$ & $0.28$ & yes\\
$0.010$ & $0.18$ & $0.56$ & yes\\
$0.020$ & $0.38$ & $1.11$ & yes\\
$0.050$ & $0.61$ & $2.76$ & yes\\
$0.100$ & $1.44$ & $5.51$ & yes\\
$0.200$ & $4.02$ & $11.05$ & no\\
\bottomrule
\end{tabular}
\end{table}

Thus raw weight distance can be anti-correlated with functional identity across a broad range of perturbation magnitudes. It only stops being backwards once the perturbation is large enough to move farther in weights than the symmetry copy.

We then ran real continued-training fine-tunes: $60$ Adam steps from the trained base model on reshuffled data. These are genuine fine-tunes, not Gaussian perturbations. The same pattern holds: fine-tunes move only about $1.1$--$1.2$ in weight distance but change function and shadow; the symmetry copy moves about $7.1$ in weights but has zero function and shadow change. Figure~\ref{fig:paradox} summarizes this outcome.

\subsection{Level-4 canonical shadows fix the spectral blind spot}

The trace/spectral shadow is sound but incomplete. We tested the failure mode directly by taking the real trained model's coupling matrix $G$ and applying a non-permutation orthogonal rotation $QGQ^T$. The spectrum is unchanged, so the spectral shadow is fooled. The Level-4 canonical coupling shadow separates the pair.

\begin{table}[h]
\centering
\caption{Level-4 canonical shadow on trained MLP couplings. The canonical shadow preserves soundness and separates a rotated coupling that the spectral shadow cannot distinguish.}
\label{tab:level4}
\begin{tabular}{lcc}
\toprule
Test & spectral distance & canonical distance\\
\midrule
Symmetry copy of model $0$ & $\approx 10^{-16}$ & $1.6\times 10^{-16}$\\
Minimum separation among $12$ trained models & $0.047$ & $0.49$\\
Orthogonally rotated real coupling & $6\times 10^{-16}$ & $0.37$\\
\bottomrule
\end{tabular}
\end{table}

In trained-like random couplings, distinct refined node signatures occurred in $2000/2000$ trials. A Weisfeiler--Lehman-style refinement was not fooled by the tested rotations in $500/500$ trials. These are not proofs of full graph-canonization completeness. They support the sharper claim: the Level-4 shadow is a generic canonical form for dense trained couplings, with the automorphic/tied-signature case isolated as the residual hard stratum.

\subsection{Identity is not alignment: a merge test that delimits scope}

A natural question is whether the canonical order also solves the harder cross-model problem of aligning two independently trained models for weight averaging. It does not, and the failure is informative. Table~\ref{tab:merge} compares loss barriers after averaging two trained models under no alignment, canonical-shadow alignment, and standard weight-matching rebasining.

\begin{table}[h]
\centering
\caption{Mean loss barrier (merge loss minus averaged endpoint loss) over independently trained model pairs. Canonical-shadow alignment barely improves on naive averaging and is far weaker than weight-matching, because it is a within-model invariant, not a cross-model correspondence.}
\label{tab:merge}
\begin{tabular}{lc}
\toprule
Alignment & mean loss barrier\\
\midrule
None (naive average) & $0.207$\\
Canonical-shadow order & $0.204$\\
Weight-matching rebasin & $0.185$\\
\bottomrule
\end{tabular}
\end{table}

Canonical order agrees with the weight-matching permutation on only $1.2\%$ of units, and canonical-matched units have weight correlation $-0.02$ against $+0.29$ for weight-matched units. The shadow is a sound within-model orbit invariant; it is not a cross-model unit correspondence. This boundary is stated as an application limit below.

\section{The shadow hierarchy}

The result is best understood as a hierarchy rather than a single metric.

\begin{center}
\begin{tabular}{lll}
\toprule
Level & Object & Role\\
\midrule
0 & Raw weights & representative-level; not invariant\\
1 & Spectral shadow & sound rejection oracle; incomplete\\
2 & Trace-power shadow & native form of Level 1\\
4 & Canonical coupling shadow & generic canonical form for $G/S_H$\\
4+ & WL-refined canonical shadow & larger-stratum canonicalization\\
\bottomrule
\end{tabular}
\end{center}

The numbering is intentional. Level 1 and Level 2 carry the same finite-dimensional information: eigenvalues are the computational readout of trace powers. Level 4 is a qualitative jump because it retains node-localized cross-terms before quotienting by order. This is the neural analogue of the move from diagonal magnitudes to phase correlations in the anyon setting.

\section{Applications}

\paragraph{Checkpoint deduplication.} If two checkpoints differ only by known gauge symmetries, raw weights can make them look unrelated. Weight shadows collapse these cases.

\paragraph{Fine-tune auditing.} Shadows can measure which blocks changed their coupling structure after continued training, alignment, distillation, or quantization.

\paragraph{Boundary test: model merging requires alignment.} It is tempting to use a Level-4 canonical order to align hidden units before weight averaging. We tested this directly and report a clean negative result, because it sharpens what the shadow is for. Aligning two independently trained microGPT models by canonical order before averaging reduced the loss barrier by only $0.003$ relative to naive averaging, whereas standard weight-matching rebasining reduced it by $0.021$ (about $7\times$ more). The canonical order agreed with the weight-matching permutation on only $1.2\%$ of units (chance level), and canonical-matched units were essentially uncorrelated ($-0.02$) while weight-matched units were clearly correlated ($+0.29$). The explanation is structural, not a tuning issue: the canonical shadow orders units by a \emph{within-model} coupling-graph signature, whereas merging requires a \emph{cross-model} correspondence between functionally matching units. These are different problems---\emph{orbit identity} (one model modulo its symmetry group) versus \emph{orbit alignment} (relating two distinct orbits). The shadow solves the first soundly; the second needs a similarity-based matcher, not an intrinsic canonical label. Merging is therefore a boundary test, not an application of this invariant.

\paragraph{Model libraries and indexing.} Blocks can be indexed by orbit-level fingerprints rather than raw tensors. This could support search over model families, variants, soups, and checkpoints.

\paragraph{Interpretability.} The coupling graph exposes hidden-unit interaction cycles, node signatures, and canonical orderings. These are structural objects, not coordinate artifacts.

\section{Limitations}

\paragraph{The modeled symmetries are explicit.} The invariant is sound for the symmetries built into it: MLP hidden-unit permutation, positive ReLU scaling, and standard head permutation. Other equivalences may exist and require other shadows.

\paragraph{Coupling completeness is not full weight-pair completeness.} The Level-4 canonical shadow is a generic canonical form for the coupling matrix $G$ modulo permutation. It is not claimed to be a complete invariant of the full pair $(F_1,F_2)$ modulo all possible functional equivalences. A stronger full-weight canonical form would retain richer descriptors such as rank-one unit operators $c_i\otimes r_i$ and their cross-correlations.

\paragraph{Tied signatures require canonicalization.} If node signatures tie, naive sorting is not enough. The correct object is weighted graph canonization, possibly with Weisfeiler--Lehman refinement or exact search inside automorphism blocks.

\paragraph{Normalization layers deserve separate treatment.} The isolated ReLU MLP scaling symmetry is exact. Full-network LayerNorm or RMSNorm can interact with scaling in ways that should be measured rather than assumed.

\paragraph{Scale remains open.} The experiments are deliberately small. Larger transformer families, residual stacks, and real checkpoint suites are the natural next tests.

\section{Conclusion}

The central lesson is simple: do not compare weights before quotienting their native symmetries. Transformer parameter space is not a plain Euclidean space of functions. It is a space of representatives with exact algebraic fibers. Raw distance can measure motion inside a fiber and mistake it for functional change.

Weight shadows repair the comparison at the right layer. The trace shadow gives a native class function: signatures before roots, cycles before coordinates. The spectral readout makes that invariant numerically measurable. The Level-4 canonical coupling shadow restores basis-sensitive cross-terms and gives a generic coordinate on the coupling quotient.

The paper's paradox now has an algebraic answer. The symmetry copy is far only as a representative. In the quotient, it is the same point. The fine-tunes are near only as raw tensors. In the shadow, they have moved. That is the point of a shadow: it measures not where the representative went, but what changed after the nuisance symmetry is removed.

\appendix

\section{Implementation sketch}

The MLP coupling shadow is computed by:
\begin{verbatim}
Gin  = fc1 @ fc1.T
Gout = fc2.T @ fc2
G    = Gin * Gout
spec = sort(eigvalsh(G))
\end{verbatim}

The Level-4 canonical coupling shadow is:
\begin{verbatim}
sig   = stack([diag(G), G.sum(axis=1), (G**2).sum(axis=1)], axis=1)
order = lexsort(sig.T[::-1])
canon = G[order][:, order].ravel()
\end{verbatim}

For tied signatures, this cheap canonicalization should be replaced by a tie-aware graph-canonization routine or a Weisfeiler--Lehman refinement followed by exact search inside ambiguous blocks.

\section{Experimental files}

The accompanying scripts are organized as follows.
\begin{itemize}[leftmargin=2em]
\item \texttt{train.py}: trains the small microGPT models from multiple seeds.
\item \texttt{experiment.py}: tests symmetry soundness and baseline spectral shadows.
\item \texttt{sweep.py}: runs perturbation magnitude sweeps and real continued-training fine-tunes.
\item \texttt{level4.py}: builds Level-4 candidate shadows and tests the cospectral/rotation failure mode.
\item \texttt{level4\_honest.py}: probes generic completeness and WL refinements.
\item \texttt{level4\_real.py}: validates Level-4 shadows on trained microGPT weights.
\end{itemize}

\begin{thebibliography}{9}

\bibitem{vaswani2017attention}
Ashish Vaswani, Noam Shazeer, Niki Parmar, Jakob Uszkoreit, Llion Jones, Aidan N. Gomez, Lukasz Kaiser, and Illia Polosukhin.
\newblock Attention Is All You Need.
\newblock \emph{Advances in Neural Information Processing Systems}, 2017.

\bibitem{karpathyMicroGPT}
Andrej Karpathy.
\newblock \emph{microgpt.py}.
\newblock Gist: \url{https://gist.github.com/karpathy/8627fe009c40f57531cb18360106ce95}.

\bibitem{ainsworth2023gitrebasin}
Samuel K. Ainsworth, Jonathan Hayase, and Siddhartha Srinivasa.
\newblock Git Re-Basin: Merging Models modulo Permutation Symmetries.
\newblock \emph{International Conference on Learning Representations}, 2023.

\bibitem{entezari2022role}
Rahim Entezari, Hanie Sedghi, Olga Saukh, and Behnam Neyshabur.
\newblock The Role of Permutation Invariance in Linear Mode Connectivity of Neural Networks.
\newblock \emph{International Conference on Learning Representations}, 2022.

\bibitem{weisfeiler1968reduction}
Boris Weisfeiler and Andrei Leman.
\newblock The reduction of a graph to canonical form and the algebra which appears therein.
\newblock 1968.

\end{thebibliography}

\end{document}
